\section{Related Work}
\label{sec:work}
\noindent \textbf{Point Cloud Networks:}
In early research, the projection-based methods~\cite{su2015multi,lawin2017deep} used to project 3D point clouds into multi-view 2D images, apply regular 2D convolutions and fuse the extracted information for 3D analysis. Alternatively, discretization-based approaches~\cite{guo2020deep} tended to convert the point clouds to voxels~\cite{huang2016point} or lattices~\cite{su2018splatnet}, and then process them using 3D convolutions or sparse tensor convolutions~\cite{choy20194d}. To avoid context loss and complex steps during data conversion, the point-based networks~\cite{qi2017pointnet,qi2017pointnet++,wang2019dynamic}  directly process point cloud data via MLP-based operations. Although current mainstream approaches in point cloud upsampling prefer utilizing MLP-related modules, in this paper, we focus on an advanced transformer structure~\cite{vaswani2017attention} in order to further enhance the point-wise dependencies between known points and benefit the generation of new points.

\noindent \textbf{Point Cloud Upsampling:}
Despite the fact that current point cloud research in low-level vision~\cite{yu2018pu,yuan2018pcn} is less active than that in high-level analysis~\cite{qi2017pointnet,hu2020randla,qi2019deep}, there exists many outstanding works that have contributed significant developments to the point cloud upsampling task. To be specific, PU-Net~\cite{yu2018pu} is a pioneering work that introduced CNNs to point cloud upsampling based on a PointNet++~\cite{qi2017pointnet++} backbone. Later, MPU~\cite{yifan2019patch} proposed a patch-based upsampling pipeline, which can flexibly upsample the point cloud patches with rich local details. In addition, PU-GAN~\cite{li2019pu} adopted the architecture of Generative Adversarial Networks~\cite{goodfellow2014generative} for the generation problem of high-resolution point clouds, while PUGeo-Net~\cite{qian2020pugeo} indicated a promising combination of discrete differential geometry and deep learning. More recently, Dis-PU~\cite{li2021point} applies disentangled refinement units to gradually generate the high-quality point clouds from coarse ones, and PU-GCN~\cite{qian2021pu} achieves good upsampling performance by using graph-based network constructions~\cite{wang2019dynamic}. Moreover, there are some papers exploring \emph{flexible-scale} point cloud upsampling via meta-learning~\cite{ye2021meta}, self-supervised learning~\cite{zhao2021sspu}, decoupling ratio with network architecture~\cite{luo2021pu}, or interpolation~\cite{qian2021deep}, \etc~As the first work leveraging transformers for point cloud upsampling, we focus on the effectiveness of PU-Transformer in performing the fundamental \emph{fixed-scale} upsampling task, and expect to inspire more future work in relevant topics. 

\noindent \textbf{Transformers in Vision:}
With the capacity in parallel processing as well as the scalability to deep networks and large datasets~\cite{khan2021transformers}, more visual transformers have achieved excellent performance on image-related tasks including either low-level~\cite{yang2020learning,chen2021pre} or high-level analysis~\cite{dosovitskiy2020image,liu2021swin,carion2020end,zhu2020deformable}. Due to the inherent gaps between 3D and 2D data, researchers introduce the variants of transformer for point cloud analysis~\cite{yew2022regtr,fan2021point,yu2022point,fan2022point,yu2021pointr}, using vector-attention~\cite{zhao2021point}, offset-attention~\cite{guo2021pct}, and grid-rasterization~\cite{mazur2021cloud}, \etc ~However, since these transformers still operate on an overall classical PointNet~\cite{qi2017pointnet} or PointNet++ architecture~\cite{qi2017pointnet++}, the improvement is relatively limited while the computational cost is too expensive for most researchers to re-implement. To simplify the model's complexity and boost its adaptability in point cloud upsampling research, we only utilize the general structure of transformer encoder~\cite{dosovitskiy2020image} to form the body of our PU-Transformer.

